% Adapted from an OpenAI Math manuscript. See PROVENANCE.md.
\section{Volume and measure consequences}\label{sec:consequences}

With \cref{thm:main} proved, we derive the additive volume inequality stated in the introduction and then apply Saroglou's transfer theorem to even log-concave measures.

\subsection{The symmetric \texorpdfstring{$L_p$}{Lp} volume inequality}
\label{subsec:lp-proof}

\begin{proof}[Proof of \cref{cor:lp}]
Put $a=|K|^{1/n}$, $b=|L|^{1/n}$, $K_0=a^{-1}K$ and $L_0=b^{-1}L$. Set
\[
 c^p=(1-\lambda)a^p+\lambda b^p,qquad
 \theta=\lambda b^p/c^p.
\]
The weighted arithmetic--geometric mean inequality yields the pointwise bound
\[
 \bigl((1-\lambda)h_K^p+\lambda h_L^p\bigr)^{1/p}
 =c\bigl((1-\theta)h_{K_0}^p+\theta h_{L_0}^p\bigr)^{1/p}
 \geq c h_{K_0}^{1-\theta}h_{L_0}^\theta.
\]
The Wulff construction is monotone in its defining function and commutes with positive dilations. By \cref{thm:main}, its value on $h_{K_0}^{1-\theta}h_{L_0}^\theta$ has volume at least $1$. The left side of \eqref{eq:lp} is therefore at least $c^p$.
\end{proof}

\subsection{Even log-concave measures and the \texorpdfstring{$(B)$}{(B)}-conjecture}
\label{subsec:sec-measure}

We now pass from volume to integration against even log-concave measures. A nonnegative Radon measure $\mu$ on $\mathbb{R}^n$ is \emph{log-concave} if
\[
 \mu((1-\theta)A+\theta B)
 \ge \mu(A)^{1-\theta}\mu(B)^\theta
 \qquad(0<\theta<1)
\]
for all compact sets $A,B\subset\mathbb{R}^n$; it is \emph{even} if $\mu(A)=\mu(-A)$ for every Borel set $A$. Radon measures are finite on compact sets, but their total mass need not be finite. We use the analogous multiplicative definition of log-concavity for nonnegative functions, so the identically zero function is included. Endpoint products in an interpolation inequality are understood as the corresponding endpoint values.

Saroglou's transfer theorem gives the logarithmic Brunn--Minkowski inequality for an even log-concave density in any dimension in which it holds for Lebesgue measure \cite[Theorem~3.1]{Saroglou2016}. Averaging the measure over small translated balls will give such a density, even when the original measure is supported on a proper subspace.

\begin{corollary}
\label{cor:measure}
Let $n\ge1$, let $\mu$ be an even log-concave Radon measure on $\mathbb{R}^n$, and let $K,L\subset\mathbb{R}^n$ be origin-symmetric convex bodies. Then, for every $0\le\lambda\le1$,
\begin{equation}\label{eq:measure-logbm}
 \mu\bigl(\mathcal{W}[h_K^{1-\lambda}h_L^\lambda]\bigr)
 \ge \mu(K)^{1-\lambda}\mu(L)^\lambda.
\end{equation}
Consequently, for every origin-symmetric convex body $K\subset\mathbb{R}^n$, the function
\begin{equation}\label{eq:b-function}
 \mathbb{R}\ni t\longmapsto \mu(\ee^tK)
\end{equation}
is log-concave.
\end{corollary}

\begin{proof}
Fix $0<\lambda<1$, and put $D=\mathcal W[h_K^{1-\lambda}h_L^\lambda]$. Let $B$ be the closed unit ball. Choose $r>0$ with $rB\subset K\cap L$. The Wulff definition then gives $rB\subset D$ as well. We may restrict $\mu$ to a sufficiently large closed ball centered at $0$ containing $2K\cup2L\cup2D$. The restriction is finite by the Radon property. It is even and log-concave because the ball is symmetric and convex, and it has the same values on $K,L,D$. The zero measure and the endpoint cases give equality, so assume the restricted measure is nonzero.

Let $\mu_\varepsilon$ be its convolution with the uniform probability on $\varepsilon B$. Its density at $x$ is $f_\varepsilon(x)=\mu(x+\varepsilon B)/|\varepsilon B|$, by Fubini and symmetry of $B$. This density is even. It is log-concave by the defining inequality for $\mu$ on compact sets, since the Minkowski interpolation of $x+\varepsilon B$ and $y+\varepsilon B$ is the translate of $\varepsilon B$ by the interpolated center. Normalizing the density to have mass one leaves the transfer inequality unchanged, since both sides are homogeneous of degree one in the measure.

Take $0<\delta<1/2$ and $\varepsilon=\delta r$. If $x\in K$ and $z\in\varepsilon B$, then $x+z\in(1+\delta)K$ by convexity and $rB\subset K$. Thus $\mu_\varepsilon((1+\delta)K)\ge\mu(K)$, and the same holds for $L$. The Wulff interpolation of $(1+\delta)K$ and $(1+\delta)L$ is $(1+\delta)D$, by homogeneity. Apply \cref{thm:main} and Saroglou's transfer to get
\[
 \mu_\varepsilon((1+\delta)D)\ge
 \mu(K)^{1-\lambda}\mu(L)^\lambda.
\]
For a point $x+z$ counted by the left side, where $z\in\varepsilon B$, one has
\[
 x\in (1+\delta)D+\varepsilon B\subset(1+2\delta)D
\]
because $rB\subset D$. Therefore the left side is at most $\mu((1+2\delta)D)$. The dilates decrease to $D$ as $\delta\downarrow0$. Continuity of the finite measure from above gives \eqref{eq:measure-logbm}.

Finally $h_{\ee^t K}=\ee^t h_K$. The Wulff interpolation of $\ee^s K$ and $\ee^t K$ is therefore $\ee^{(1-\lambda)s+\lambda t}K$. Apply \eqref{eq:measure-logbm} to this pair to obtain \eqref{eq:b-function}.
\end{proof}
